\section{Related Work}
\label{sec:related-work}

% =========================================================================
% Four threads plus a self-improvement aside and a position summary. The
% hierarchical-RL background was set up in Section \ref{sec:background-hrl};
% this section positions our contribution against each thread.
% =========================================================================

We position the contribution against four adjacent lines of work:
hierarchical reinforcement learning, recursive and tree-structured
neural network encoders, Monte Carlo Tree Search applied to language
model reasoning, and verifier-anchored proof systems that employ
structural decomposition. In each case the design choice we develop
(implicit recursion via parent-pointer references in a flat state
representation, with a depth-limit prior selecting the productive
regime) departs from the dominant architectural pattern in that line,
and the departure is itself the contribution.

\subsection{Hierarchical Reinforcement Learning}
\label{sec:rw-hrl}

The hierarchical reinforcement learning literature reaches for hierarchy
as an \emph{architectural} solution: separate sub-policies, value
decompositions with explicit cross-subtask coupling, or two-level
manager-worker splits. Section~\ref{sec:background-hrl} summarized three
representative threads; here we make the positional contrast explicit.

\paragraph{Options framework.}
The options framework of Sutton, Precup, and Singh
\citep{sutton1999between} introduces temporally extended actions, each
with its own internal policy and termination condition. An MDP equipped
with options becomes a semi-Markov decision process whose Bellman
equations operate over option-completion intervals; the option-critic
architecture \citep{bacon2017option} learns options end-to-end. Our
verbs are not options. A verb such as $\textsc{Decompose}$ executes in a
single decision step and does not enter a sub-policy with its own
termination condition; the action set $\Aset$ stays seven-element
categorical throughout. The analog of an option in our setting is the
recursive subtree that $\textsc{Decompose}$ induces, and that subtree is
planned over by MCTS acting as the policy improvement operator over
$\Qpsi$, not by a separately parameterized sub-policy. Option-style
temporal abstraction is obtained from the policy improvement operator's
search machinery, not from new architectural components in the value
function or behavior policy.

\paragraph{MAXQ value decomposition.}
Dietterich's MAXQ decomposition \citep{dietterich2000hierarchical}
factors the value function along a fixed task hierarchy, exposing shared
subtasks for transfer and faster planning. Our setting admits an
analogous decomposition along the implicit hypothesis tree
(Section~\ref{sec:background-hrl}), but we do not impose it
architecturally; $\Qpsi$ is a multi-head value function over the fixed
seven-verb action set and sees only the flat serialized state. The MAXQ
framing is useful as an analytical lens on sample efficiency (a
productive recursion in our sense is precisely one whose MAXQ-style
completion can be assembled from verifier-anchored leaf signals), but
the architecture does not depend on the decomposition being explicit.
The MCTS chance-node statistics aggregate over the trajectory
naturally, without requiring the value function to factor along the
tree.

\paragraph{Feudal and goal-conditioned hierarchies.}
Feudal networks \citep{vezhnevets2017feudal} separate a manager
proposing goals from a worker achieving them, with the manager
operating at a coarser timescale; goal-conditioned hierarchical RL
\citep{nachum2018data} learns goal embeddings end-to-end. Unlike feudal
networks, our manager (the policy choosing $\textsc{Decompose}$) and
worker (the policy choosing $\textsc{ComputeSubexpr}$ or
$\textsc{InvokeVerifier}$ on a decomposed sub-claim) share the same
frozen LLM substrate, operate at the same per-step timescale, and are
both expressed as choices over $\Aset$. The hierarchy lives in the
representation of the state, not in the parameterization of the policy.

\paragraph{Position.}
Across all three threads, hierarchy is introduced as an additional
architectural ingredient. We deliberately avoid this. Implicit hierarchy
via parent-pointer identifier references in a flat state, with the
categorical action set $\Aset$ unchanged, is methodologically more
parsimonious: no new sub-policy parameters, no separate sub-policy
selection step, no architectural commitment to a particular
decomposition. The price is that the encoder (the frozen LLM) must
recover the tree structure from serialized depth and parent
annotations; the design assumes that modern decoder-only LLMs'
in-context attention is up to this task.
Section~\ref{sec:method-serialization} develops that argument and
Section~\ref{sec:experimental-plan}'s ablation grid provides its
falsification.

\subsection{Recursive Neural Networks and Tree Encoders}
\label{sec:rw-trees}

A separate tradition addresses tree-structured input directly at the
encoder level. Recursive neural networks \citep{socher2013recursive}
compose hidden states along a parsed structure, applying a learned cell
recursively from the leaves to the root; Tree-LSTMs
\citep{tai2015improved} extend the construction with LSTM gating to
mitigate vanishing-gradient issues over deep trees. Graph attention
networks \citep{velickovic2018gat} generalize to arbitrary
parent-pointer graphs via message passing, producing a representation
explicitly aware of graph topology. Set transformers
\citep{lee2019set} attack a related problem, providing
permutation-invariant attention over unordered collections. Each
architecture builds a representational prior into its parameter
sharing: respect tree composition, respect graph topology, respect
permutation symmetry.

Our position is deliberately contrary: we avoid all of these. The
frozen LLM, operating on flat token serialization with depth markers
and parent-pointer identifier annotations, is our tree encoder. We
argue in Section~\ref{sec:method-serialization} that in-context causal
self-attention is a sufficient substitute given the modest depth
($\depthlim$ on the order of three) and modest token counts our
deliberation states reach. The reason this is plausible is that the
base LLM's pretraining already exposes it to deeply structured text:
HTML and XML with nested tags, source code with nested scopes,
mathematical notation with subscripts and parenthesized sub-expressions,
JSON documents with parent-child references. The in-context retrieval
behavior that handles those constructions also handles the explicit
\texttt{parent=} and \texttt{depth=} annotations in our serialization
template, with no specialized architectural machinery. The advantages of
the flat-token approach are architectural simplicity (no extra encoder,
no representation drift from a separate encoder family, no fine-tuning
required to align encoder and generator) and a single hidden state on
which the heads ($\Qpsi$, $\ptheta$, and the verb-instantiator) all read
off.

\subsection{MCTS for Language Model Reasoning}
\label{sec:rw-mcts-llm}

The broader test-time-compute scaling literature
\citep{snell2024scaling, openai2024o1} has shown that deployed
reasoning systems can substantially improve accuracy by spending
inference compute on deliberate generation, with chain-of-thought
prompting \citep{wei2022cot} as the prototypical baseline.
MCTS-style search over LLM continuations is one specific way to spend
that compute. Feng et al.\ \citep{feng2023alphazerollm} explore
AlphaZero-style tree search as a means of guiding LLM decoding and
training, treating partial generations as MCTS nodes and using a value
network trained alongside the policy. Xie et al.\ \citep{xie2024monte}
use MCTS to improve reasoning via iterative preference learning, with
search trajectories furnishing pairwise preferences for downstream
fine-tuning. Tree of Thoughts \citep{yao2023tree} introduces tree
search over intermediate reasoning steps and is the closest prior work
in structural ambition: thoughts are explored in a branching tree
rather than a linear chain. The closest published system to the
framework paper's design is rStar-Math
\citep{guan2025rstarmath}, which combines small base language models
with MCTS-driven self-training and a process reward model trained from
MCTS rollouts: the same family of small-LM, search-guided,
distillation-from-search approaches that our framework inhabits.

Our work is distinct in three ways. First, we operate within an
approximate $Q$-learning framework where MCTS is the policy improvement
operator, not the substrate of deliberation itself. The framework paper
positions $\Qpsi$ as the central learned object and MCTS as one
sample-efficient procedure for producing better Bellman targets;
deliberation proceeds by repeated value-and-action choice over the
categorical verb set $\Aset$, with MCTS supplying $\Qmcts$ and
$\pimcts$ as training targets back to $\Qpsi$ and $\ptheta$. The
LLM-MCTS systems above tend to invert this relationship, treating MCTS
or the equivalent tree search as the deliberation mechanism itself and
fitting evaluator networks to support it. Second, prior LLM-MCTS work
tends to treat the state as a token sequence (one MCTS node per partial
generation) rather than as a structured object that can itself be
tree-recursive. Tree of Thoughts uses tree-shaped exploration over
reasoning paths but does not exercise structural decomposition
\emph{within} a single state, where a hypothesis splits into
sub-hypotheses that are themselves entities in the state. rStar-Math,
similarly, operates over token-level or step-level decoding with
search nodes keyed to partial generations of natural-language
reasoning; the recursive structure available to its MCTS is the
generation tree of the next reasoning step, not an implicit hypothesis
tree carried in a structured state. The tree in those systems is the
MCTS tree; the tree we discuss is implicit in the state structure and
lives at every MCTS node, so depth limits, identifier-based action
targeting, and the productive-versus-unproductive distinction apply at
the deliberation-state level rather than at the decoding-step level.
Third, we treat structural recursion as an explicit design dimension
with a depth-limit prior. Prior LLM-MCTS systems, including rStar-Math
and Tree of Thoughts, either disallow recursion in the deliberation
state (the state is a flat reasoning prefix) or permit it implicitly
without analyzing when recursion is productive and when it degenerates
into rumination. The present paper makes that distinction
load-bearing.

\subsection{Recursive Decomposition in Formal Proof Systems}
\label{sec:rw-proof}

The verifier-anchored proof-search literature offers the closest
operational analogs to our setting. AlphaProof
\citep{deepmind2024alphaproof} achieved silver-medal performance on
International Mathematical Olympiad problems by combining a Lean-based
formal verifier with tree search over tactic applications; Lean's
tactic structure imposes a natural recursive decomposition on the
proof, where each tactic may open subgoals proved recursively. LeanDojo
\citep{yang2023leandojo} provides the infrastructure on which much of
this work runs: a proof-state environment that exposes Lean's
intermediate goals to a retrieval-augmented language model, with proof
construction proceeding through tactic application and subgoal
resolution. HyperTree Proof Search
\citep{lample2022hypertree} is the most direct prior work on
bounded-depth recursive subgoal trees: a hypertree representation
allows each proof step to fan out into multiple subgoal conjunctions
proved independently, and search proceeds with an explicit budget on
hypertree depth and a learned value function over open subgoals. The
mechanism is the closest published analog to our depth-limit-as-prior
treatment of recursion.

Our position relative to this line is methodological rather than
competitive. AlphaProof and LeanDojo operate over the Lean tactic
hierarchy, inheriting recursion from Lean's type system: the structural
decomposition is given by the proof assistant. HyperTree Proof Search
makes the recursive subgoal structure explicit as a hypertree data
structure carried by the search algorithm, with bounded-depth
hypertree expansion as the search-time control. We make the recursion
mechanism itself explicit, as a design dimension of the deliberation
MDP independent of any particular target domain. The state-level
contrast against HyperTree is that our parent pointers carry the tree
implicitly inside a flat structured state schema rather than as an
explicit hypertree DAG built by the planner; the search-level
contrast is that bounded depth in our setting is a mask on the
MDP's recursive verbs rather than on the planner's expansion budget,
so the productive-versus-unproductive distinction is enforced at the
environment level before any specific MCTS configuration is fixed. Our parent-pointer
references in a structured state schema, our depth-limit prior, and our
distinction between productive and unproductive recursion
(Definition~\ref{def:productive}) are intended as the methodological
framework behind which recursive-decomposition systems implicitly
operate. The implication is not that our framework substitutes for
Lean. The implication is that systems like AlphaProof exhibit the
productive-recursion regime when their tactic trees are well-conditioned
(each subgoal is verifier-checkable and contributes terminal proof
reward to the trajectory) and degenerate toward the unproductive regime
when the tactic search produces deep chains of meta-tactic application
without leaf-level verifier signal. Our contribution is to expose that
distinction operationally and to propose the depth limit as a mechanism
for selecting the productive mode within a finite sample budget, in a
way that transfers across target domains where the verifier is not a
formal proof assistant.

\subsection{Self-Improvement and Compositional Generalization}
\label{sec:rw-self}

Self-improvement frameworks such as STaR \citep{zelikman2022star}
bootstrap reasoning by generating candidate rationales, filtering by
verifier outcome, and training on the survivors. Larger-scale
verifier-anchored RL pipelines such as DeepSeek-R1
\citep{guo2025deepseekr1} extend the same principle to rule-based
verifiers for math and code without a learned reward model,
deliberately avoiding the Goodhart attack surface that attends
preference-trained critics. The expert-iteration view of the framework
paper, in which MCTS-improved rollouts produce targets for $\Qpsi$ and
$\ptheta$, can be read as a search-augmented version of this loop. The productive-recursion mechanism we develop is
one structural ingredient by which such a system might learn to compose:
STaR's rationales are flat chains; ours are tree-structured by
construction, with each node addressing a verifier-checkable
sub-problem. Whether the tree structure improves compositional
generalization across problems (reusing a verified sub-claim or a
decomposition strategy across problem families) is an open question we
flag but do not test here.

\subsection{Position Summary}
\label{sec:rw-summary}

The four threads above each touch the present contribution at a
different angle. Hierarchical RL shares the structural-decomposition
intuition but commits to explicit architectural hierarchy; we provide
implicit hierarchy via the state representation alone. The
tree-structured-encoder tradition shares the recognition that nested
inputs deserve specialized treatment but commits to encoder-side
recursion; we let the frozen LLM's in-context attention substitute.
MCTS-for-LLM work shares the use of tree search but treats MCTS as the
substrate of deliberation; we treat it as the policy improvement
operator over $\Qpsi$ in the expert-iteration sense established by the
framework paper. Verifier-anchored proof systems share the operational
shape of recursive subgoal resolution but inherit recursion from a
specific formal target (Lean); we extract the recursion mechanism as a
domain-independent design dimension. The present paper is small in
scope, developing one design choice from the framework paper in depth.
Its novelty is the explicit productive-versus-unproductive distinction
and the depth-limit prior that operationalizes it within an approximate
$Q$-learning system over an LLM-induced context MDP.
